\scanpage{50}
\mspage{6}
\uline{Remarque} En fait, on a un foncteur plus précis, vers une catégorie de graphes munis d'une structure où la partie $\rho$ de la pondération est \deleted{\unsure{complétée}} \inserted{remplacée} par une structure plus précise, à savoir \deleted{une \unsure{application}} \inserted{une \deleted{\unsure{famille}}} famille
\deleted{\diagram{\begin{tikzcd}[ampersand replacement=\&,row sep=small,column sep=small]
A\,(\widetilde R,\alpha,\unsure{p|\rho}:R\arrow[r]\&\vec A\,)\arrow[d,equal]\\
\&\vec A/\sigma
\end{tikzcd}}}
\[
 (R_a)_{a\in A}
\]
\margin{$\omega(a)\mathrel{\mathop=^{\rm def}}\{\vec a\in\vec A/\vec a\text{ au dessus de }a\}$}
d'ensembles, et pour $\forall a\in A$, si $\omega(a)$ désigne l'\unsure{ens|un} des
\deleted{\unsure{donné d'une permutation de}}
``orientations'' de $a$ ($\operatorname{card}\omega(a)\in\{1,2\}$), une application
\deleted{$R$, et $\unsure{p|\rho}$ une application $R\longrightarrow A$ \unsure{conj. à}}
\[
 \omega\longmapsto u_\omega:\quad \omega(a)\longrightarrow\operatorname{Circ}(R_a)
\]
ens des permutations circulaires sur $R_a$
\deleted{et \unsure{d'autre} \unsure{une} \unsure{\illegible{1} une application injective|application injective}}
\[
 \deleted{R/\alpha\lhook\joinrel\longrightarrow\vec A}
\]
avec la condition que si $\omega,\omega'\in\omega(a)$, $\omega\ne\omega'$, alors
\[
 \deleted{u_\omega u_{\omega'}=\operatorname{id}_{\unsure{R_a|R}}}
 \qquad u_{\omega'}=u_\omega^{-1}
\]

On peut aussi considérer que la donnée des suppléments revient à la donnée de
\[
 \begin{array}{r@{}c@{}l}
 &r&:A\longrightarrow\mathbb N\\[-.3ex]
 &\Vert&\\[-.3ex]
 a\longmapsto&\operatorname{Card}&\unsure{p|\rho}^{-1}(a)
 \end{array}
\]
(la ``pondération-arêtes''), et la donnée d'une famille de torseurs
\[
 (R_a)_{a\in A},\quad \deleted{\unsure{\omega(a)}}
\]
$R_a$ un torseur sous $\mathbb Z/r(a)\mathbb Z\mathbin{\wedge}^{\omega(a)}_{\{\pm1\}}$
On associe à $(X,\widetilde\Sigma,R)$ une telle structure supplémentaire :
\[
 R_a=R\cap\widetilde\Sigma_a
\]
($\widetilde\Sigma_a$ la comp. connexe de $\widetilde\Sigma$ \deleted{\unsure{marquée}} \inserted{qui correspond à} $a\in A\simeq\pi_0(\widetilde\Sigma)$),
\deleted{$\unsure{R=R(\Sigma C)|R=\pi_0(\widetilde\Sigma)}$,}
\deleted{\diagram{\begin{tikzcd}[ampersand replacement=\&,row sep=small,column sep=large]
R\arrow[r]\arrow[d,equal]\&S\arrow[d,equal]\\
\unsure{\Gamma}\arrow[r]\&\pi_0(\widetilde\Sigma)
\end{tikzcd}}}
\deleted{\unsure{la donnée de} $\widetilde\Sigma$}
$\omega\longmapsto u_\omega$ défini de façon évidente.
\deleted{et $\unsure{p|\rho}$ déduit de l'ordre circulaire \unsure{dans|sur} les $R_a$}

\uline{NB} Même si $r(\omega)=1$ ou $r(\omega)=2$, le fait de disposer
\mspage{49}
